\section{从辽东到寿春：边境的胜利进入中枢}

\subsection{237--245：扩展的半径与继承的裂缝}

二三七年，公孙渊拒绝魏命、自立
燕王。次年司马懿围襄平，
公孙氏覆灭；同在二三八年，
倭女王卑弥呼遣使至魏，魏朝廷授予册封
和物品。前者显示魏可以在东北投入远征资源，后者显示洛阳的名号可以沿海路
抵达更远的政治关系；两者都不等于一张边境以外的控制地图。有关倭地、航线和
岛屿的对应，必须保留传世记载与后出区域史之间的距离。

远征成功并未解决洛阳的继承问题。二三九年
曹叡去世，幼帝曹芳即位，曹爽与司马懿
同受遗诏。此后的辅政不是皇帝旁边的一项职位，而是通向中军、诏令和任官的
接口。正始年间所立的三体
石经，让经书校写在洛阳太学留下可见痕迹；残石足以说明官学与公开校读的
存在，不能推出全国学术整齐一致。

孙吴的裂缝来自另一种继承。二四一年
太子孙登去世，二四二年
孙和为太子、孙霸封鲁王，两个成年皇子
各自吸引宾客与官僚。陆逊后来围绕此事的谏争，在二四五年他
去世后更成为后世叙述的焦点；能够确定的
是继承竞争已把朝廷的关系带到江防资源旁边，而不是每一句传下来的劝谏都能
还原为当日的会议纪录。

魏也没有把边境优势自由转化为进攻。二四四年曹爽
攻蜀至兴势而撤，王平等据险的
防守使山道再度限制中枢的雄心。毌丘俭二四四至二四五年的
高句丽之役攻破丸都城，东川王避走；
战报可确认攻城与撤离，却不应将此写成东北地方从此被魏稳定整合。

\subsection{247--253：辅政者如何变成政权的入口}

蜀汉在二四七年由姜维趁陇西羌胡起事出兵，
重新把北伐重心推向陇右。这不是
诸葛亮路线的简单复写：羌胡起事给了机会，郭淮等魏军的反应以及粮道仍限定了
机会的大小。

二四九年，司马懿发动高平陵政变，
控制洛阳并罢黜曹爽。政变的场所、近畿军队和辅政名义缺一不可；它既不是抽象
的“世族胜利”，也不只是宫门中的一次偶然。它把魏的中枢接口改交司马氏，
并使淮南、关中和东北的将领不得不重新判断谁能发令、谁能保护他们。

江东的选择更为急促。二五〇年孙权
废孙和、赐死孙霸，改立孙亮，
以清洗结束二宫之争。二五一年
司马懿死，司马师继承中军与辅政地位；
王凌随即谋立楚王曹彪，事泄后
投降自尽，曹彪被迫死。两处
继承危机都说明，制度的连续性不能只靠遗诏或血缘，必须有人维持军队与地方的
服从。

二五二年孙权去世，孙亮继位、诸葛恪
辅政。诸葛恪旋即在东兴堤击退魏军，利用
水陆地形保住濡须方向；战胜给新辅政者带来威望，却也把更大的军事期待压到
建业。二五三年他围合肥新城，
因疫病、补给和守备撤退，同年孙峻
杀诸葛恪，接掌禁军和朝政。从东兴到
建业的转换并非两段无关故事：前线的损耗会把对辅政者的判断带回宫门。

\subsection{254--260：淮南和建业的反复夺权}

魏的皇帝名义继续缩小。二五四年司马师
废曹芳、立曹髦；翌年毌丘俭、文钦
在淮南起兵，失败后文钦等投吴。
司马师在镇压后病死，
司马昭接掌辅政与中军。淮南不是魏、
吴之间一条简单的国界线：当地军镇、洛阳的政变和吴援军的选择在这里碰撞。

吴的禁军同样在重组。二五六年
孙峻死，孙綝继掌朝政；二五七年诸葛诞
据寿春反司马昭并求吴援。二五八年，
孙綝废孙亮、立孙休；寿春同时
陷落，诸葛诞被杀，吴援军受挫。吴
能够把军队送到淮南，却不能因此替叛军维持一座被魏围困的城。

二六〇年曹髦率近侍出宫，试图讨伐司马昭，
途中被杀，曹奂继立。此事让魏的名义
秩序暴露出最直接的限度：皇帝仍可行动，却不能在关键时刻调动足以改变中枢
控制权的资源。三国的下一阶段，将从这种已被改写的北方中枢出发。

\begin{sourcenote}[同时代的边界不是现代国界]
辽东、高句丽、倭和淮南的条目主要由《三国志》及裴注、后出的区域史和考古
研究互照。朝贡、册封、攻城或一次驻军可证实接触和行动，不能自动证明长期、
均质的统治。正始石经和东兴的叙述也须分别按残石、传记和编年材料所能承担的
尺度来读。
\end{sourcenote}
